\documentclass[10pt, graphics, aspectratio=169, table]{beamer}
\usepackage{hyperref}
\usepackage{qrcode}
\usepackage[normalem]{ulem}
% -- Change to this year's color! --
\definecolor{ese}{RGB}{255, 176, 0}


\usetheme{metropolis}
\setbeamercolor{frametitle}{bg=ese}

\title{Open Source Development}
\author{Eric Wolf}
\date{ESE \the\year{}}
\institute{Nerd::101 - ESE - IFSR - TU Dresden}
\titlegraphic{\hfill\includegraphics[height=2cm]{../logo}}

\begin{document}
    \maketitle

    \begin{frame}{Outline}
        \tableofcontents
    \end{frame}

    \section{Introduction}
    \begin{frame}{What is Open Source?}
        \begin{columns}
            \column{0.7\textwidth}
                \begin{itemize}
                    \item practice of publishing software publicly alongside its source code
                    \item central goals: allowing use, study, modification and distribution
                    \item different implementations via software licenses
                \end{itemize}
            \column{0.3\textwidth}
                \centering
                \includegraphics[height=0.25\textheight]{img/firefox.png}
                \center\tiny\url{https://www.firefox.com}
                \includegraphics[width=\textwidth]{img/libreoffice.png}
                \center\tiny\url{https://www.libreoffice.org}
                \includegraphics[height=0.25\textheight]{img/notepadpp.png}
                \center\tiny\url{https://notepad-plus-plus.org}
        \end{columns}
    \end{frame}

    \begin{frame}{Why should you care?}
        Open Source \dots
        \begin{itemize}
            \item \dots allows for easy collaborative development
            \item \dots allows the community to take control if necessary (via forking)
            \item \dots makes it hard to implement unwanted `features' or security backdoors
        \end{itemize}

        \textbf{Open Source software empowers its users: you!}
    \end{frame}

    \section{Commonly used licenses}
    \begin{frame}{Permissive licenses}
        \begin{columns}
            \column{0.6\textwidth}
                \begin{itemize}
                    \item allow almost anything, including commercial use and relicensing
                    \item require copies to include a copyright and permission notice
                    \item do not provide any warranty or liability
                    \item examples: MIT License and Apache License 2.0
                \end{itemize}
            \column{0.4\textwidth}
                \includegraphics[width=\textwidth]{img/apache.png}
                \center\tiny\href{https://commons.wikimedia.org/wiki/File:ASF_Logo.svg}{Apache Software Foundation Logo}
        \end{columns}
    \end{frame}

    \begin{frame}{The problem with permissive licenses}
        \begin{itemize}
            \item freedoms granted by permissive licenses do not have to be preserved
            \item often summarized as use, study, share and improve
            \item such software can be redistributed without granting them to its users
            \item copyleft licenses ensure they are preserved by redistribution
        \end{itemize}

        \textbf{Copyleft licenses are sometimes referred to as `viral' licenses}
    \end{frame}

    \begin{frame}{Copyleft licenses}
        \begin{columns}
            \column{0.6\textwidth}
                \begin{itemize}
                    \item similar to permissive licenses, but more strict
                    \item require the source code to be disclosed on distribution and the declaration of changes
                    \item require modifications and larger works to be released under a compatible license
                    \item examples: GNU General Public License and its variants
                \end{itemize}
            \column{0.4\textwidth}
                \includegraphics[width=\textwidth]{img/GPL.png}
                \center\tiny\href{https://commons.wikimedia.org/wiki/File:GPLv3_Logo.svg}{GNU General Public License Logo}
        \end{columns}
    \end{frame}

    \section{Example: GitHub}
    \begin{frame}{What is GitHub?}
        \begin{columns}
            \column{0.7\textwidth}
            \begin{itemize}
                \item platform for collaborative development using git
                \item widely used in the open source community
                \item includes issue trackers, website hosting, CI/CD, \dots
                \item is not open source itself (proprietary)
            \end{itemize}
            \column{0.3\textwidth}
                \includegraphics[width=\textwidth]{img/github.png}
                \center\tiny\url{https://github.com}
        \end{columns}
    \end{frame}

    \begin{frame}{Collaborative development on GitHub}
        Contribute changes to another repository by \dots
        \begin{enumerate}
            \item copy the original repository by forking it
            \item clone and push your changes to the forked repository
            \item create a pull request to the original repository
            \item wait for owner of original repository to accept the pull request
        \end{enumerate}

        \textbf{The actual process depends on the project you want to contribute to}
    \end{frame}

    \begin{frame}{Showtime}
        \textbf{A little demonstration \dots}
    \end{frame}

    \section{What to do next}
    \begin{frame}{What to do next}
        \begin{columns}
            \column{0.6\textwidth}
                \begin{itemize}
                    \item look into what open source software you are using (think about transitive dependencies)
                    \item participate and be active in their communities
                    \item support open source projects (via code, work or money)
                \end{itemize}

            \column{0.4\textwidth}
                \centering
                \includegraphics[height=0.75\textheight]{img/dependency.png}
                \center\tiny\url{https://xkcd.com/2347}
        \end{columns}
    \end{frame}

    \begin{frame}{Choose a license}
        \center\qrcode[height=2in]{https://choosealicense.com}
        \center\url{https://choosealicense.com}
    \end{frame}
\end{document}
